\section{Portfolio choice under the certificate}\label{sec:portfolio-choice}

The preceding section establishes how much portfolio variance the
observed information geometry can certify under the maintained
transmission model.
The investor's task is now to turn that risk statement into an
allocation rule without treating the bound as an estimated covariance matrix.
At the benchmark with $\delta=0$, the natural decision is to choose the
feasible capital allocation with the smallest certified upper bound in
\eqref{eq:p3-raw-cap}.

Let $\mathcal W$ be a nonempty compact subset of the long-only simplex,
possibly incorporating upper-weight, sector, or turnover constraints.
Each $x\in\mathcal W$ is a vector of capital weights, whereas $q(x)$
records the corresponding normalized shares of the
perfect-positive-dependence volatility budget defined in
\eqref{eq:p3-risk-weights}.
Define the certified-variance objective
\begin{equation}\label{eq:p3-certified-objective}
  \operatorname{CV}(x)
  :=A(x)^2\left\{1-\mathcal C(q(x))\right\}.
\end{equation}
The information-certified portfolio solves
\begin{equation}\label{eq:p3-certified-portfolio}
  x^{\star}\in
  \operatorname*{arg\,min}_{x\in\mathcal W}\operatorname{CV}(x).
\end{equation}
Equation \eqref{eq:p3-certified-portfolio} is a minimum-risk problem.
Expected returns do not enter its objective, constraints, or identifying
argument.
The objective is economically appropriate because it minimizes a
certified risk cap rather than a point estimate of otherwise unknown
portfolio variance.

The two components of \eqref{eq:p3-certified-objective} make the
allocation trade-off explicit.
The term $A(x)^2$ is the variance benchmark under perfect positive
dependence, so reducing $A(x)$ lowers the same marginal-volatility
benchmark used by an inverse-volatility portfolio.
The deduction $A(x)^2\mathcal C(q(x))$ rewards allocations whose
information distributions certify that their latent risks cannot all be aligned.
Inverse volatility is therefore the cleanest empirical benchmark: it
uses the first channel but not the second.

The first mathematical requirement is that this decision rule
actually has a solution.

\begin{theorem}[Existence of a certified-variance minimizer]
  \label{thm:p3-certified-minimizer}
  Every nonempty compact feasible subset $\mathcal W$ of the long-only
  simplex admits a minimizer of the continuous objective
  \eqref{eq:p3-certified-objective}, and the pointwise variance certificate
  holds at that minimizer.
\end{theorem}

For fixed distance and marginal-scale inputs, $\operatorname{CV}$ is continuous
on $\mathcal W$.
The set is nonempty and compact, so Weierstrass gives a minimizer.
The pointwise certificate in \eqref{eq:p3-raw-cap} then holds at that
minimizer.

Existence makes the rule well defined on any compact feasible set;
curvature is a separate guarantee about how the standardized problem
can be solved.
That curvature is a property of the observed geometry rather than an
assumption about returns.
For a symmetric zero-diagonal matrix $d$, conditional negative definiteness
means
\[
  \sum_i\sum_j u_i\,u_j\,d_{ij}^2\le0
  \quad\text{whenever }\sum_i u_i=0.
\]

\begin{theorem}[Certificate convexity]\label{thm:p3-certificate-convexity}
  If the pairwise floor matrix $\ell$ is conditionally negative definite, then
  $q\mapsto1-\mathcal C(q)$ is convex along mixtures of normalized risk weights.
\end{theorem}

For every tangent direction with $\1^\top u=0$,
\[
  \frac{d^2}{dt^2}\mathcal C(q+tu)=u^\top\ell^2u\le0.
\]
Hence $\mathcal C$ is concave and $1-\mathcal C$ is convex on the simplex.

By Schoenberg's criterion the condition holds exactly when the double-centred
matrix
$-\tfrac12(I-\tfrac1n\1\1^{\top})\ell^2
(I-\tfrac1n\1\1^{\top})$
is positive semidefinite, equivalently when the floors embed isometrically in a
Hilbert space.
It is therefore checkable directly from the reported W2 matrix, without
returns and without an estimated parameter, and it holds for the frozen matrix
used here.
Under this condition, minimizing the standardized certificate over a convex
feasible set is a convex program.

A third guarantee concerns sensitivity to the carrier scale rather
than existence or curvature.
In the zero-slack boundary case, the carrier constant drops out of
the normalized decision.

\begin{corollary}[Carrier invariance of the normalized allocation]
  \label{cor:p3-carrier-invariance}
  Let $\tau_i=0$ for every asset. Then $\ell_{ij}=L^{-1}W_2(C_i,C_j)$ and
  \begin{equation}\label{eq:p3-zero-slack-certificate}
    \mathcal C(q)=\frac{1}{2L^2}\,q^{\top}W_2^2\,q,
  \end{equation}
  so for any $L>0$ and any feasible set $\mathcal Q$ of normalized risk
  weights,
  \[
    \operatorname*{arg\,max}_{q\in\mathcal Q}\mathcal C(q)
    =\operatorname*{arg\,max}_{q\in\mathcal Q}q^{\top}W_2^2q.
  \]
\end{corollary}

The factor $L^{-2}$ is positive, so scaling the objective does not change its
maximizers.
Under zero slack, $L$ determines how much variance the certificate deducts
but not which normalized risk allocation achieves the largest deduction.
That allocation is determined by the observed W2 geometry.

The carrier-invariance statement concerns normalized risk weights
$q$, which are shares of the perfect-dependence volatility budget,
not the capital weights $x$ held by the investor.
The raw objective retains $L$ through the interaction between $A(x)^2$ and
$\mathcal C(q(x))$, so converting an invariant risk allocation into capital
still requires the marginal scales $\sigma_i$.
When slack is nonzero, $\tau_i+\tau_j$ enters the floor additively and the
truncation depends jointly on $L$ and $\tau$.
Together, the existence, curvature, and carrier-invariance results
turn the certificate into an implementable decision rule while
keeping its dependence on the maintained inputs explicit.
\Cref{sec:empirical-design} next constructs the zero-slack allocation
from observed news geometry, and \Cref{sec:news-variance-benchmark}
evaluates its realized variance against conventional portfolio benchmarks.
